\section{Contribuciones, uso de inteligencia artificial y reproducibilidad}

\subsection{Contribuciones}
La tabla resume el registro final de \texttt{docs/contributions.md}. Cada aporte se
vincula con archivos, pruebas, resultados o materiales verificables del repositorio.

\begin{table}[ht]
\centering\small
\caption{Registro de contribuciones.}
\label{tab:contribuciones}
\begin{tabularx}{\linewidth}{@{}l X@{}}
\toprule
Integrante & Contribución registrada \\
\midrule
Pablo Daniel Barillas Moreno & Coordinación e integración; definición del problema; investigación del algoritmo genético; implementación inicial del simulador, la CLI y la optimización restringida; reproducibilidad y cierre \\
Jorge Palacios & Pipeline reproducible de NHAMCS 2022 (descarga, preparación, calibración y pruebas), agregados en \texttt{data/processed/}, \texttt{docs/05\_data\_calibration.md} y entradas NHAMCS de la bibliografía \\
Adrián Ricardo González Muralles & Correcciones del simulador (cola tras calentamiento, drenaje, capacidad efectiva), diseño factorial, candidata congelada, 180 ejecuciones, validación frente a NHAMCS, figuras y pruebas \\
Andrés Rafael Chivalán & Interfaz Streamlit en cinco páginas, formularios de ejecución controlada, visualizaciones, evidencia precomputada, descargas y pruebas de interfaz \\
Javier Andrés Chen & Informe técnico en LaTeX, generación reproducible de tablas y figuras (\texttt{Output/report/}, \texttt{Output/scripts/build\_report\_assets.py}), presentación y guion cronometrado \\
\bottomrule
\end{tabularx}
\end{table}

\subsection{Uso de inteligencia artificial}
Se declara en \texttt{docs/ai\_usage.md}: OpenAI Codex apoyó la planificación, la estructura
del repositorio, la revisión de código, las correcciones del simulador, el pipeline
experimental y la documentación inicial; Claude Code apoyó la documentación del pipeline de
calibración de NHAMCS. Para este informe se usó además Claude (Anthropic) como apoyo para
estructurar y redactar el borrador y para escribir el script que convierte la evidencia en
tablas y figuras. Las cifras proceden exclusivamente de \texttt{evidence/final/} y de
\texttt{data/processed/}, y los integrantes son responsables de revisarlas y de poder
explicarlas. Las herramientas de IA no se usaron como fuente de datos ni como autoridad
bibliográfica.

\subsection{Reproducibilidad del informe}
Las tablas, las macros con cifras y las figuras de este informe se regeneran con
\begin{center}\small\texttt{python Output/scripts/build\_report\_assets.py}\end{center}
y el documento se compila con \texttt{latexmk -xelatex main.tex} desde
\texttt{Output/report/}.
Ninguna cifra de las secciones de resultados se escribió a mano.
